\scanpage{263}

\mspage{134}

$Z\xrightarrow[\sim]{u}T_H\hookrightarrow T_G\xrightarrow{\mathrm{can}}\widetilde Z$. Il faudra de plus, \unsure{enfin|encore}, décrire l'hom. $c:\widetilde Z\longrightarrow Z$.

Rappelons-nous que l'on a

\[\left\{\begin{aligned}
T_0&=\left\{g\in\unsure{\mathfrak{S}_S|\mathscr{G}_S}\ \middle|\ 
\begin{aligned}&g\text{ normalise }H\subset\mathfrak S_{\unsure{S|z}}\\&g\text{ centralise } Z\subset\mathfrak S_S\\&g\deleted{\unsure{\in}}s_0=s_0\\&gt_0\in Z.t_0\end{aligned}\right\}\deleted{\Longleftrightarrow}\\
T_1&=\left\{g\in\mathfrak S_S\ \middle|\ 
\begin{aligned}&g\text{ normalise }H\\&g\text{ centralise }Z\\&gt_0=t_0\\&gs_0\in Z.s_0\end{aligned}\right\}
\end{aligned}\right.\ .\tag*{(137)}\]

\margin{$T_G=\left\{g\in\mathfrak S_S\ \middle|\ \begin{smallmatrix}g\text{ norm. }H\\g\text{ centr. }Z\\gs_0\in Zs_0\\gt_0\in Zt_0\end{smallmatrix}\right.$}

Or pour un espace homogène $S$ sous un groupe $H\inserted{\subset\mathfrak S_S}$, $S$ \unsure{ponctué|pointé} par une ``origine $s_0$'', les groupe des $g\in\mathfrak S_S$ qui normalisent $H$ et fixent $s_0$, i.e. des autom. de la situation

\[(S,H\subset\mathfrak S_S,s_0\in S)\longleftrightarrow(H,U_0=H_{s_0}),\quad\text{est can.}\]

isomorphe au groupe des automorphismes \inserted{$\bar g\deleted{\unsure{g}}$} de $H$ qui \deleted{\unsure{stabilisent}} normalisent $U_0$. \deleted{\unsure{Le ca}} Pour un \deleted{tel} $g$, dire qu'il \deleted{centralise $Z\subset\mathfrak S_S$, défini} \inserted{normalise un sous-groupe} $Z\subset\mathfrak S_S$ commutant à \deleted{\unsure{$G$}} $H$, défini par un sous-groupe $\underset{\deleted{\unsure{T_H}}}{\underline{B_{0H}}}\supset U_0$ du normalisateur de $U_0$, \unsure{se} signifie que l'autom. correspondant $\bar g$ de $H$ normalise $B_{0H}$. -- et l'autom \deleted{\unsure{induit sur}} $\inserted{\operatorname{int}(g)|Z}$ s'identifie alors à $\operatorname{int}(\bar g)|(B_{0H}/U_0)$.

Donc dire que $g$ \uline{centralise} $Z$, signifie que $\bar g$ normalise $B_{0H}$ et induit l'identité sur $B_{0H}/U_0$.

Il reste à exprimer la condition $gt_0\in Z.t_0$ en termes de $\bar g$, or \deleted{$gt_0$} $t_0$ correspond à $\sigma.U_0$, $gt_0$ à $\bar g(\sigma).U_0$, et \deleted{on voit qu'il existe} \inserted{$\lambda t_0$ ($\lambda\in Z$) \unsure{corres}}

% archive sheet number: 262
